\subsection*{Assignment 4 -- Data Warehouse Schema Design}

L’Assignment 4 richiedeva il passaggio da un database operativo a un
data warehouse orientato all’analisi, con l’obiettivo di progettare
uno schema multidimensionale in grado di analizzare il numero di stream
registrati a un mese dalla pubblicazione di una canzone
(\textit{Streams1Month}).

Lo schema fornito nella consegna era indicato come riferimento non vincolante,
lasciando libertà progettuale nella definizione della granularità,
delle dimensioni e delle relazioni. Seguendo l’approccio discusso a lezione,
la progettazione è partita dall’analisi delle \textit{business questions}
e delle query che si intendevano supportare.

Da tale analisi è emerso che la misura centrale del data warehouse è unica:
\begin{itemize}
    \item \textbf{Streams1Month}, ovvero il numero di stream registrati a un mese
    dalla pubblicazione del brano.
\end{itemize}

A partire da questa misura sono state individuate le dimensioni rilevanti
per l’analisi:
\begin{itemize}
    \item \textbf{DimSong}, che rappresenta la canzone pubblicata;
    \item \textbf{DimArtist}, che rappresenta gli artisti che hanno partecipato
    alla canzone;
    \item \textbf{DimDate}, che rappresenta la data di pubblicazione del brano;
    \item \textbf{DimArtistGeography}, che rappresenta il luogo di nascita
    dell’artista;
    \item \textbf{DimLyrics}, che rappresenta il testo del brano e le sue componenti;
    \item \textbf{DimSymphony}, che rappresenta le caratteristiche acustiche
    del brano;
    \item \textbf{DimAlbum}, che rappresenta il contesto editoriale della canzone,
    non necessariamente coincidente con un album tradizionale.
\end{itemize}

Un’ulteriore scelta progettuale rilevante è stata quella di non utilizzare
le \textit{natural key} presenti nei dataset originali, in quanto fortemente
duplicate e poco affidabili. Sono state pertanto introdotte chiavi surrogate
univoche, generate tramite UUID, al fine di garantire coerenza, stabilità
e corretta identificazione delle dimensioni.

\subsubsection*{Evoluzione dello schema e scelte progettuali}

In una prima fase progettuale è stata considerata una fact table denominata
\textit{PublishedSong}, con granularità definita come
\emph{una riga per canzone pubblicata da un artista}, collegata alle dimensioni
Song, Artist, Date, Lyrics e Symphony.

In tale schema, la gestione degli artisti in featuring sarebbe stata possibile
esclusivamente tramite l’introduzione di una \textit{bridge table} tra
\textit{DimSong} e \textit{DimArtist}. Tuttavia, questo approccio presentava
alcune criticità rilevanti. In primo luogo, l’uso della bridge table avrebbe
condotto a una struttura che non aderiva pienamente né allo \textit{star schema}
né allo \textit{snowflake schema}, aumentando il numero di join necessari per
l’esecuzione delle query analitiche e riducendo l’efficienza complessiva del
sistema.

In secondo luogo, la gestione dei featuring risultava meno naturale dal punto
di vista analitico, poiché il ruolo dell’artista (artista primario o featuring)
non era direttamente associato alla misura analizzata, ma derivato da una
struttura intermedia.

Sono state inoltre considerate altre alternative progettuali, tra cui:
\begin{itemize}
    \item l’introduzione di un attributo \textit{Role} nella bridge table;
    \item l’utilizzo di una \textit{degenerate dimension} per distinguere tra
    artisti primari e featuring;
    \item la creazione di una dimensione di gruppo di artisti tra la fact table
    e la dimensione Artist.
\end{itemize}

Sebbene concettualmente corrette, tali soluzioni introducevano una maggiore
complessità strutturale o rendevano meno immediate le analisi per artista,
risultando meno adatte agli obiettivi dell’assignment.

\subsubsection*{Schema finale del Data Warehouse}

La soluzione finale adottata è stata la creazione di una fact table denominata
\textit{FactParticipation}, che modella esplicitamente la partecipazione di un
artista a una canzone.

La granularità finale è definita come:
\begin{quote}
\emph{una riga = partecipazione di un artista a una canzone in una determinata data}.
\end{quote}

Di conseguenza, la stessa canzone può comparire più volte all’interno della fact
table, una volta per ciascun artista coinvolto. Questo comportamento è atteso
ed è il risultato di una granularità più fine, scelta in modo consapevole per
rappresentare correttamente la relazione molti-a-molti tra artisti e canzoni,
evitando duplicazioni semantiche delle misure.

Il ruolo dell’artista è rappresentato tramite l’attributo \textit{IsPrimary},
modellato come \textit{degenerate dimension} all’interno della fact table:
\begin{itemize}
    \item \textit{IsPrimary = 1} identifica l’artista principale;
    \item \textit{IsPrimary = 0} identifica gli artisti in featuring.
\end{itemize}

La fact table è collegata direttamente alle dimensioni \textit{DimSong},
\textit{DimArtist} e \textit{DimDate}, mentre le dimensioni Song e Artist
generano strutture a \textit{snowflake schema} rispettivamente verso
\textit{DimAlbum}, \textit{DimLyrics}, \textit{DimSymphony} e
\textit{DimArtistGeography}. L’insieme delle foreign key identifica univocamente
i record della fact table, in accordo con quanto discusso a lezione.

La dimensione \textit{DimDate} è stata collegata direttamente alla fact table,
in quanto rappresenta il riferimento temporale dell’evento analizzato
(stream a un mese dalla pubblicazione). Altre date semanticamente differenti,
come la data di rilascio dell’album, sono state mantenute come attributi
all’interno delle rispettive dimensioni, evitando ambiguità e join aggiuntivi.

\subsubsection*{Implementazione}

Una volta definito lo schema multidimensionale del Data Warehouse, è stata
realizzata anche la sua implementazione fisica sul database server utilizzando
SQL Server Management Studio. In particolare, è stato scritto ed eseguito uno
script SQL per la creazione delle tabelle dimensionali e della fact table,
in piena coerenza con lo schema logico descritto.

Le tabelle sono state inizialmente create vuote, definendo esplicitamente:
\begin{itemize}
    \item gli attributi e i relativi tipi di dato;
    \item le chiavi primarie;
    \item le chiavi esterne che implementano le relazioni tra fact table e dimensioni.
\end{itemize}

Le foreign key riflettono correttamente la struttura del Data Warehouse,
garantendo allineamento tra modello concettuale, logico e implementazione fisica.
Sono stati inoltre introdotti vincoli e scelte sui tipi di dato per garantire
coerenza e qualità dei dati. In particolare:
\begin{itemize}
    \item l’attributo \textit{Gender} della dimensione \textit{DimArtist} è stato
    vincolato ai soli valori ammessi (\textit{M} o \textit{F}) tramite una
    \textit{check constraint};
    \item i tipi di dato numerici sono stati scelti in base alla natura degli
    attributi, utilizzando \textit{INT} per conteggi discreti,
    \textit{FLOAT} per misure continue e \textit{BIT} per attributi booleani
    come \textit{IsPrimary};
    \item la misura \textit{Streams1Month} è stata definita come attributo numerico
    adeguato a rappresentare valori di conteggio di grandi dimensioni.
\end{itemize}

Questa fase ha permesso di separare in modo chiaro la progettazione dello schema
dalla fase di popolamento dei dati, che viene affrontata negli assignment
successivi. Il caricamento effettivo dei dati nel Data Warehouse è infatti
realizzato nell’Assignment~6, tramite strumenti ETL, a partire dai file CSV
prodotti nell’Assignment~5.
